The implementation was successfully validated using multiple demonstration scripts.
The primary objective of these tests was to ascertain whether the framework could seamlessly spawn, initialize, and control both drones and cars simultaneously within the exact same Unreal Engine 5 instance, without triggering the core engine conflicts that plagued the baseline single-domain Cosys-AirSim.\newline\newline
The test \texttt{\seqsplit{multi\_agent\_drone\_car}} instantiates two independent client objects from the Cosys-AirSim Python API. The extension allows these clients to bind to dedicated Localhost API server ports: port 41451 is exclusively assigned to handle the multirotor control signals, while port 41452 was bound to the car commands.\newline
During the active execution phase, both vehicles achieve real-time, asynchronous operation. The script commands \texttt{Drone1} to arm, execute an automated takeoff sequence (\texttt{takeoffAsync}), and perform trajectory tracking via \texttt{moveToPositionAsync}. In parallel, \texttt{Car1} is simultaneously supplied with continuous \texttt{CarControls} inputs (throttle set to 0.5, steering set to 0.5), initiating a curved forward movement. \newline\newline
The test \texttt{car\_hit\_drone} specifically validates the physical interaction and collision detection between heterogeneous agents. In this scenario, \texttt{Drone1} is placed in a stationary state on the ground, acting as a passive obstacle. \texttt{Car1} is then commanded to move forward at full throttle directly towards the UAV. The results demonstrate that the Unreal Engine physics engine, as extended by the framework, correctly processes the collision: the car physically impacts the drone and reflects off it, confirming that the separate API domains do not isolate the agents from the global physics environment.\newline\newline
The test \texttt{cars\_intersection\_crash} evaluates multi-vehicle interaction within the ground domain. Three cars are instantiated: two navigating along a horizontal axis and one along a vertical axis, creating a four-way intersection scenario. The script commands all vehicles to accelerate simultaneously towards the center. Upon impact, the framework accurately simulates the transfer of momentum and torque, resulting in realistic physical behaviors such as vehicle spinning and lateral displacement post-collision, further validating the robustness of the multi-agent physics integration.\newline\newline
The scenario \texttt{drone\_hit\_car} provides a dual-domain collision test where the UAV acts as the kinetic agent. The script initializes \texttt{Car1} in a braking state and calculates its vertical position to allow \texttt{Drone1} to lock its altitude precisely to the car's chassis level. Upon impact, the framework demonstrates consistent collision geometry: the drone is physically repelled by the car's collider, reflecting the impact. While the stationary car's displacement is negligible due to the significant mass differential and simulated tire friction, the test confirms that the UAV's flight controller and the global physics engine successfully register the collision event across the heterogeneous multi-agent setup.\newline\newline
The test \texttt{drones\_rapid\_cross} focuses on high-velocity multi-agent coordination within the aerial domain. Three UAVs are programmed to perform a synchronized, high-speed crossing maneuver at the same altitude. During the execution, the system maintains stable control over all three drones until a collision occurs between two of them. The results show that the framework correctly detects the high-speed impact; the involved drones become physically interlocked, demonstrating that the physics engine accurately handles complex mesh-to-mesh interactions even at significant relative velocities.\newline\newline
The \texttt{full\_chaos} test serves as a stress and stability benchmark for the multi-port API extension. This scenario instantiates a large number of heterogeneous agents, both cars and drones, and commands them to move in divergent, semi-random directions. The simulation results confirm that the framework can simultaneously handle multiple physics-driven interactions: some vehicles successfully navigate open spaces while others collide with static environment geometry (e.g., buildings, walls). Throughout this high-load phase, the communication between the distributed API servers (multirotor and car) remains synchronized, with no observed latency spikes or command drops, validating the scalability of the proposed extension.\newline\newline
The \texttt{many\_vehicles} scenario validates the system's ability to manage a coordinated fleet of heterogeneous agents. In this test, the framework controls two UAVs and two UGVs simultaneously, while other inert vehicles are present in the environment map. The script executes a synchronized escort pattern where both drones and cars advance in a forward trajectory. The telemetry logs demonstrate that the framework maintains independent and precise control over each specific vehicle instance (e.g., \texttt{Drone1}, \texttt{Drone2}, \texttt{Car1}, \texttt{Car2}) across their respective API ports, ensuring that the presence of multiple active agents does not lead to command crosstalk or resource exhaustion within the Unreal Engine instance.\newline\newline
The final test, \texttt{swarm\_takeoff}, examines the synchronization of multi-agent aerial maneuvers. In this scenario, three drones are commanded to perform a simultaneous takeoff, followed by a rapid vertical ascent to an altitude of 15 meters. The results confirm that the framework can broadcast and synchronize identical high-level commands across multiple UAV instances without jitter or significant timing offset, validating the system's capability for swarm-based research.\newline\newline
The results are definitively positive. Polling the \texttt{getCarState} and \texttt{getMultirotorState} functions confirm that the vehicles achieve stable initialization and control. Both the UAV and UGV successfully coexist and navigate their distinct physical domains, thereby thoroughly fulfilling the primary objective of the extension. While the framework demonstrates high functional reliability across all tested scenarios, future work will focus on optimizing long-duration stability under complex collision environments to ensure maximum simulation uptime.